%% Simplification and tuning of common-slide-1.tex
\synctex=1
\documentclass[xcolor=dvipsnames]{beamer}
\mode<presentation>



\usepackage[font=itshape]{quoting}
\usepackage[]{graphicx}
\usepackage[dvipsnames]{xcolor}

\usepackage{colortbl}

%%\usepackage{CSIC3}

\usetheme{default}
% \selectcolormodel{gray} %% Nope, no red for ACTIVITIES, etc
\usecolortheme{dove}
\setbeamercolor{frametitle}{fg=black}



\setbeamercolor{frametitle}{bg=white!90}
\setbeamertemplate{frametitle}{%
  \vspace*{6.3pt} %% increase/decrease for space above slide title
  %% In the next, wd controls width of text
  \begin{beamercolorbox}[wd=1.05\textwidth, dp=0.024cm]{frametitle}
    \usebeamerfont{frametitle}\insertframetitle
    \noindent
    \vspace*{-13pt}
    \rule{1.05\textwidth}{1.2pt}
  \end{beamercolorbox}
  % \vspace*{0.1cm}
}


\setbeamertemplate{headline}{
   \begin{beamercolorbox}[ignorebg,ht=1.87ex,dp=0.0ex]{section in head/foot}
     \colorbox{gray!10}{
       \insertsectionnavigationhorizontal{\paperwidth}{}{}
     }
   \end{beamercolorbox}
}

%% In headline, current section in bold, non-current non-bold
\setbeamerfont{section in head/foot}{series=\bfseries}
\setbeamertemplate{section in head/foot shaded}{\mdseries\color{gray}\usebeamertemplate{section in head/foot}}


\setbeamertemplate{footline}{
  \leavevmode%
  \hbox{%
\begin{beamercolorbox}[wd=1.00\paperwidth,ht=2.25ex,dp=1ex,right]{subsection in head/foot}%
      \usebeamerfont{subsection in  head/foot}\insertframenumber{} / \inserttotalframenumber\hspace*{2ex}
    \end{beamercolorbox}}%
  \vskip0pt%
}

%% The headline with the section title. It also
%% moves page numbers to the right, relative to CSIC3. OK, fine.
%% Actually I do not need it, as it is all above: the headline and footline.
% \useoutertheme[subsection=false]{mysmoothbars3}


%%%%% or the next three and maybe the following three for navigation balls
% \usetheme{myFrankfurt}
% \useoutertheme[subsection=false]{mysmoothbars}
% \usecolortheme{whale} %% wolverine was original
%% whale OK
\makeatletter
  \beamer@compressfalse
\makeatother
%%%%% Issues with multiple lines, etc for dots
%%%% http://tex.stackexchange.com/questions/35796/how-can-i-get-multiple-lines-of-frame-dots-in-beamer-navigation




\setbeamercolor{itemize item}{fg=black}
%\setbeamertemplate{itemize item}[circle] %% [ball] ball does not change to black
% \setbeamertemplate{itemize item}[ball]
%% Larger circle
\setbeamertemplate{itemize item}{\Huge\raise-5.05pt\hbox{\textbullet}}


%% None of the back, forward, etc
\setbeamertemplate{navigation symbols}{}


\usepackage{array}
%% trying to fix interaction with pfg if using notes on second screen
\usepackage[absolute,overlay]{textpos}
% \usepackage[overlay]{textpos}

%\usepackage[spanish]{babel} % can do funny things with \lim \min
\usepackage[english]{babel}
\usepackage[latin1]{inputenc}
\usepackage{times}
\usepackage[T1]{fontenc}
\usepackage{setspace}
% Or whatever. Note that the encoding and the font should match. If T1
% does not look nice, try deleting the line with the fontenc.
%%\usepackage{fancybox}

\usepackage{gitinfo}

\usepackage{url}

\usepackage{bm}
\newcommand{\cyan}[1]{{\textcolor {cyan} {#1}}}
\newcommand{\blu}[1]{{\textcolor {blue} {#1}}}
\newcommand{\Burl}[1]{\blu{\url{#1}}}
\newcommand{\red}[1]{{\textcolor {red} {#1}}}
\newcommand{\green}[1]{{\textcolor {green} {#1}}}
\newcommand{\mg}[1]{{\textcolor {magenta} {#1}}}
\newcommand{\og}[1]{{\textcolor {PineGreen} {#1}}}
\newcommand{\code}[1]{\texttt{\slshape\footnotesize #1}}
\newcommand{\Code}[1]{\texttt{\slshape #1}}
\newcommand{\myverb}[1]{{\footnotesize\texttt {\textbf{#1}}}}
\newcommand{\Rnl}{\ +\qquad\ }
\newcommand{\Emph}[1]{\emph{\mg{#1}}}
\newcommand{\R}{{\sf R}}

\newenvironment{myitemize}
{\begin{itemize}}{\end{itemize}}

\newcommand{\mydes}[2][]{
  \vspace*{-1pt}
%%  \renewcommand{\baselinestretch}{0.7}\small\normalsize
  \begin{description}\item[{#2}]{#1}\end{description}
\vspace*{-8pt}
}



\usepackage{tikz}
\usetikzlibrary{arrows,shapes,positioning}

\tikzset{
  % Define standard arrow tip
  >=stealth',
  % Define style for boxes
  punkt/.style={
    rectangle,
    rounded corners,
    draw=black, %very thick,
    %maximum width=6em,
    %text width=6.5em,
    minimum height=2em,
    text centered},
  punkt2/.style={
    rectangle,
    rounded corners,
    draw=black, %very thick,
    text width=2.35cm,
    %text width=6.5em,
    minimum height=2em,
    inner sep = 2pt,
    text centered},
  punkt3/.style={
    rectangle,
%    rounded corners,
    draw=black, %very thick,
%    text width=1.cm,
    %text width=6.5em,
    minimum height=2em,
    inner sep = 3pt,
    text centered},
  punkk/.style={
    rectangle,
    rounded corners=3mm,
    draw=black, %very thick,
    %text width=6.5em,
    minimum height=2em,
    text centered},
  ell1/.style={
    %ellipse,
    %draw = black, very thick,
    text width=1.7cm,
    inner sep = 0pt,
    text centered},
  % Define arrow style
  pil/.style={
    ->,
    thick,
    shorten <=2pt,
    shorten >=2pt,}
}


%% Increase spacing between items in itemize and enumerate
\AddToHook{cmd/itemize/after}{\addtolength{\itemsep}{1.2\baselineskip}}
\AddToHook{cmd/enumerate/after}{\addtolength{\itemsep}{1.2\baselineskip}}

%% Increase nested item spacing
\setbeamertemplate{itemize/enumerate subbody begin}{\vspace{.75\baselineskip}}

%% First item
\setbeamertemplate{itemize/enumerate body begin}{\vspace{1.0\baselineskip}}
%% Last
\setbeamertemplate{itemize/enumerate body end}{\vspace{1.5\baselineskip}}

%% When you need to modify spacing
%%  \begin{itemize}\addtolength{\itemsep}{-1.8ex}

% %%    Nested lists
% \AddToHook{cmd/itemize/after}{\addtolength{\topsep}{2.5\baselineskip}}
% \AddToHook{cmd/enumerate/after}{\addtolength{\topsep}{1.5\baselineskip}}

%% Though maybe this? to limit it just to nested lists
% \makeatletter
% \AddToHook{cmd/itemize/after}{%
%   \ifnum\@listdepth>1\relax
%     \addtolength{\topsep}{1.5\baselineskip}%
%   \fi
% }
% \makeatother



%% for fbox
\setlength{\fboxrule}{2pt} %% thicker
\setlength{\fboxsep}{0pt} %% no padding



% maxwidth is the original width if it is less than linewidth
% otherwise use linewidth (to make sure the graphics do not exceed the margin)
\makeatletter
\def\maxwidth{ %
  \ifdim\Gin@nat@width>\linewidth
    \linewidth
  \else
    \Gin@nat@width
  \fi
}
\makeatother


%%%%%%%%%%%%%%%%%%%%%%%%%%%%%% knitr
%%% Some slides in -II and probably -III still have some knitr stuff
%%% It could be easily removed (just formats some text).
\usepackage{alltt}

\definecolor{fgcolor}{rgb}{0.345, 0.345, 0.345}
\newcommand{\hlnum}[1]{\textcolor[rgb]{0.686,0.059,0.569}{#1}}%
\newcommand{\hlstr}[1]{\textcolor[rgb]{0.192,0.494,0.8}{#1}}%
\newcommand{\hlcom}[1]{\textcolor[rgb]{0.678,0.584,0.686}{\textit{#1}}}%
\newcommand{\hlopt}[1]{\textcolor[rgb]{0,0,0}{#1}}%
\newcommand{\hlstd}[1]{\textcolor[rgb]{0.345,0.345,0.345}{#1}}%
\newcommand{\hlkwa}[1]{\textcolor[rgb]{0.161,0.373,0.58}{\textbf{#1}}}%
\newcommand{\hlkwb}[1]{\textcolor[rgb]{0.69,0.353,0.396}{#1}}%
\newcommand{\hlkwc}[1]{\textcolor[rgb]{0.333,0.667,0.333}{#1}}%
\newcommand{\hlkwd}[1]{\textcolor[rgb]{0.737,0.353,0.396}{\textbf{#1}}}%
\let\hlipl\hlkwb

\usepackage{framed}
\makeatletter
\newenvironment{kframe}{%
 \def\at@end@of@kframe{}%
 \ifinner\ifhmode%
  \def\at@end@of@kframe{\end{minipage}}%
  \begin{minipage}{\columnwidth}%
 \fi\fi%
 \def\FrameCommand##1{\hskip\@totalleftmargin \hskip-\fboxsep
 \colorbox{shadecolor}{##1}\hskip-\fboxsep
     % There is no \\@totalrightmargin, so:
     \hskip-\linewidth \hskip-\@totalleftmargin \hskip\columnwidth}%
 \MakeFramed {\advance\hsize-\width
   \@totalleftmargin\z@ \linewidth\hsize
   \@setminipage}}%
 {\par\unskip\endMakeFramed%
 \at@end@of@kframe}
\makeatother

\definecolor{shadecolor}{rgb}{.97, .97, .97}
\definecolor{messagecolor}{rgb}{0, 0, 0}
\definecolor{warningcolor}{rgb}{1, 0, 1}
\definecolor{errorcolor}{rgb}{1, 0, 0}
\newenvironment{knitrout}{}{} % an empty environment to be redefined in TeX

%%%%%%%%%%%%%%%%%%%%%%%%%%%%%%%%%%%%%%%% end knitr



\author[Ram�n D�az-Uriarte]{Ram�n D�az-Uriarte}
\institute[]{Dept. Bioqu�mica\\
Universidad Aut�noma de Madrid\\
Madrid, Spain\\
\texttt{r.diaz@uam.es}\\
\Burl{https://ligarto.org/rdiaz}
}





%%% Local Variables:
%%% mode: latex
%%% ispell-local-dictionary: "en_US"
%%% coding: iso-8859-15
%%% End:
